% Part: first-order-logic
% Chapter: axiomatic-deduction
% Section: proving-things

\documentclass[../../../include/open-logic-section]{subfiles}

\begin{document}

\iftag{FOL}
      {\olfileid{fol}{axd}{pro}}
      {\olfileid{pl}{axd}{pro}}

\olsection{نمونه‌هایی از \usetoken{P}{derivation}}


\begin{ex}
فرض کنید می‌خواهیم $(\lnot !D \lor !E) \lif (!D \lif !E)$ را اثبات
کنیم. روشن است که این فرمول نمونهٔ هیچ‌یک از اصول موضوع ما نیست؛ پس باید
برای !!{derive} کردن آن از قاعدهٔ \MP{} استفاده کنیم. تنها قاعده‌ای که
اینجا نیاز داریم MP است؛ این قاعده با داشتن $!A$ و $!A \lif !B$ به ما
اجازه می‌دهد~$!B$ را توجیه کنیم. یک راهبرد آن است که
\olref[prp]{ax:lor3} را با $!A$ به‌جای $\lnot !D$، $!B$ به‌جای $!E$ و
$!C$ به‌جای $!D \lif !E$ به کار ببریم؛ یعنی نمونهٔ
\[
(\lnot !D \lif (!D \lif !E)) \lif ((!E \lif (!D \lif !E)) \lif ((\lnot
!D \lor !E) \lif (!D \lif !E))).
\]
چرا؟ دو بار کاربرد MP بخش پایانی را به دست می‌دهد، که همان مطلوب ماست.
به‌آسانی می‌بینیم $\lnot !D \lif (!D \lif !E)$ نمونه‌ای از
\olref[prp]{ax:lnot2} و $!E \lif (!D \lif !E)$ نمونه‌ای از
\olref[prp]{ax:lif1} است. پس !!{derivation} ما چنین است:
\begin{derivation}
  1. &  $\lnot !D \lif (!D \lif !E)$ & \olref[prp]{ax:lnot2} \\
  2. & $(\lnot !D \lif (!D \lif !E)) \lif {}$\\
  &\qquad $((!E \lif (!D \lif !E)) \lif ((\lnot !D \lor !E) \lif (!D \lif !E)))$ & \olref[prp]{ax:lor3}\\
  3. & $(!E \lif (!D \lif !E)) \lif ((\lnot !D \lor !E) \lif (!D \lif !E))$ & 1, 2, \MP\\
  4. & $!E \lif (!D \lif !E)$ & \olref[prp]{ax:lif1}\\
  5. & $(\lnot !D \lor !E) \lif (!D \lif !E)$ & 3, 4, \MP
\end{derivation}
\end{ex}

\begin{ex}\ollabel{ex:identity}
بیایید بکوشیم !!a{derivation} برای $!D \lif !D$ بیابیم. این فرمول
نمونهٔ یک اصل موضوع نیست؛ پس باید با \MP{} آن را !!{derive} کنیم.
\olref[prp]{ax:lif1} اصل موضوعی به‌شکل~$!A \lif !B$ است که می‌توانیم
~\MP را بر آن اعمال کنیم. البته برای سودمندبودن این کار، $!B$ای که
\MP{} در این حالت به‌عنوان گام درست توجیه می‌کند باید~$!D \lif !D$
باشد، زیرا می‌خواهیم همین را !!{derive} کنیم. این بدان معناست که $!A$
نیز باید $!D$ باشد؛ یعنی می‌توانیم نمونهٔ زیر از
\olref[prp]{ax:lif1} را بررسی کنیم:
\[
!D \lif (!D \lif !D)
\]
برای اعمال \MP، باید مقدمهٔ دوم متناظر، یعنی~$!A$، را نیز توجیه کنیم.
اما در مورد ما این مقدمه~$!D$ خواهد بود و نمی‌توانیم~$!D$ را به‌تنهایی
!!{derive} کنیم. پس به راهبرد دیگری نیاز داریم.

اصل موضوع دیگر که فقط~$\lif$ را دربر دارد
\olref[prp]{ax:lif2} است، یعنی
\[
(!A \lif (!B \lif !C)) \lif ((!A \lif !B) \lif (!A \lif !C))
\]
می‌توانیم با دو بار کاربرد \MP{} به آخرین شرطیِ تودرتو برسیم. باز این
بدان معناست که نمونه‌ای از \olref[prp]{ax:lif2} می‌خواهیم که در آن
$!A \lif !C$ همان $!D \lif !D$، یعنی !!{formula} مورد نظر ما، باشد.
پس البته $!A$ و $!C$ هر دو~$!D$ هستند. $!B$ را چگونه برگزینیم تا هم
$!A \lif (!B \lif !C)$ و هم $!A \lif !B$، یعنی در مورد ما
$!D \lif (!B \lif !D)$ و $!D \lif !B$، نیز !!{derivable} باشند؟
عبارت نخست، مستقل از اینکه $!B$ را چه بگیریم، از پیش نمونه‌ای از
\olref[prp]{ax:lif1} است. و $!D \lif !B$ نیز نمونهٔ دیگری از
\olref[prp]{ax:lif1} خواهد بود، اگر $!B$ را $(!D \lif !D)$ بگیریم. پس
!!{derivation} ما چنین است:
\begin{derivation}
  1. & $!D \lif ((!D \lif !D) \lif !D)$ & \olref[prp]{ax:lif1}\\
  2. & $(!D \lif ((!D \lif !D) \lif !D)) \lif {}$\\
  & \qquad $((!D \lif (!D \lif !D)) \lif (!D \lif !D))$ & \olref[prp]{ax:lif2}\\
  3. & $(!D \lif (!D \lif !D)) \lif (!D \lif !D)$ & 1, 2, \MP\\
  4. & $!D \lif (!D \lif !D)$ & \olref[prp]{ax:lif1}\\
  5. & $!D \lif !D$ & 3, 4, \MP
\end{derivation}
\end{ex}

\begin{ex}\ollabel{ex:chain}
گاهی می‌خواهیم نشان دهیم از چند !!{formula} دیگر در~$\Gamma$،
!!a{derivation} برای !!a{formula} وجود دارد. برای نمونه، نشان می‌دهیم
می‌توان $!A \lif !C$ را از $\Gamma = \{!A \lif !B, !B \lif !C\}$
!!{derive} کرد.
\begin{derivation}
  1. & $!A \lif !B$ & \Hyp\\
  2. & $!B \lif !C$ & \Hyp\\
  3. & $(!B \lif !C) \lif (!A \lif (!B \lif !C))$ & \olref[prp]{ax:lif1} \\
  4. & $!A \lif (!B \lif !C)$ & 2, 3, \MP\\
  5. & $(!A \lif (!B \lif !C)) \lif {}$\\
  & \qquad $((!A \lif !B) \lif (!A \lif !C))$ & \olref[prp]{ax:lif2}\\
  6. & $((!A \lif !B) \lif (!A \lif !C))$ & 4, 5, \MP\\
  7. & $!A \lif !C$ & 1, 6, \MP
\end{derivation}
سطرهایی که برچسب «\Hyp» (مخفف «فرض») دارند نشان می‌دهند !!{formula}
آن سطر !!a{element}~$\Gamma$ است.
\end{ex}

\begin{prop}
  \ollabel{prop:chain} اگر $\Gamma \Proves !A \lif !B$ و
  $\Gamma \Proves !B \lif !C$، آنگاه $\Gamma \Proves !A \lif !C$.
\end{prop}

\begin{proof}
  فرض کنید $\Gamma \Proves !A \lif !B$ و $\Gamma \Proves !B \lif
  !C$. آنگاه !!a{derivation} برای $!A \lif !B$ از~$\Gamma$ و نیز
  !!a{derivation} برای~$!B \lif !C$ از~$\Gamma$ وجود دارد. این دو را
  با الحاقشان در یک !!{derivation} واحد ترکیب کنید. اکنون سطرهای ۳--۷
  !!{derivation} نمونهٔ پیشین را بیفزایید. این !!a{derivation} برای
  $!A \lif !C$---که آخرین سطر !!{derivation} جدید است---از~$\Gamma$
  است. توجه کنید اگر ارجاع به شمارهٔ سطر~۱ با ارجاع به آخرین سطر
  !!{derivation} برای~$!A \lif !B$، و ارجاع به شمارهٔ سطر~۲ با ارجاع به
  آخرین سطر !!{derivation} برای~$!B \lif !C$ جایگزین شود، توجیه‌های
  سطرهای ۴ و~۷ همچنان معتبر می‌مانند.
\end{proof}

\begin{prob}
با ارائهٔ !!{derivation}s از اصول موضوع نشان دهید احکام زیر برقرارند:
\begin{enumerate}
\item $(!A \land !B) \lif (!B \land !A)$
\item $((!A \land !B) \lif !C) \lif (!A \lif (!B \lif !C))$
\item $\lnot(!A \lor !B) \lif \lnot !A$.
\end{enumerate}
\end{prob}



\end{document}
